Un échantillon de \qty{\v{Ve}}{mL} d'eau de piscine contient \qty{\v{mc}}{mg} de chlore. On considère que \qty{1}{L} d'eau a une masse de \qty{1}{kg}, de sorte que \qty{1}{ppm} correspond à \qty{1}{mg/L}. La piscine contient \qty{\v{Vp}}{m^3} d'eau ($\qty{1}{m^3} = \qty{1000}{L}$).

\begin{enumerate}[label=\alph*)]
    \item Calculez la concentration du chlore dans la piscine, en ppm.
    \item Quelle masse totale de chlore, en grammes, la piscine contient-elle ?
    \item La concentration recommandée est de \qty{3}{ppm}. Quelle masse de chlore faut-il ajouter pour l'atteindre ?
\end{enumerate}
